
\section{从函数集合的概率描述开始}
高斯过程可以理解为对未知函数的一组联合概率假设：在任意有限输入位置上，函数值向量服从多元高斯分布。均值函数说明先验中心，核函数说明两个输入处函数值的相关结构。这里的“先验”是模型假设，不等于已经知道真实函数；数据会通过条件分布改变预测。

设训练观测满足 $\mathbf y=\mathbf f+\boldsymbol\epsilon$，独立高斯噪声方差为 $\sigma_n^2$，函数先验均值为0、协方差矩阵为 $K$。对新输入 $x_*$，与训练位置的协方差向量为 $\mathbf k_*$。固定核与噪声参数后，条件均值和潜在函数方差为
\[
\mu_* = \mathbf k_*^\mathsf T(K+\sigma_n^2I)^{-1}\mathbf y,
\qquad
v_* = k(x_*,x_*)-\mathbf k_*^\mathsf T(K+\sigma_n^2I)^{-1}\mathbf k_*.
\]
实际计算应通过线性方程求解或Cholesky分解，不显式求逆。若预测下一次带噪观测，还需在适当独立假设下加上观测噪声方差。核参数若由数据估计，简单条件方差通常没有完整积分其估计不确定性。

\section{两个观测点的一次完整计算}
取输入0、1，观测0、1，RBF核 $k(x,z)=\exp[-(x-z)^2/2]$，噪声标准差0.1。则
\[
K+\sigma_n^2I=\begin{pmatrix}1.01&e^{-1/2}\\e^{-1/2}&1.01\end{pmatrix},
\quad\mathbf k_*=(e^{-1/8},e^{-1/8})^\mathsf T
\]
对应新输入0.5。数值求解得到预测均值约0.5459，潜在函数方差约0.03645；预测一个带同等独立噪声的新观测时，方差还要增加0.01。均值没有恰好等于两点算术平均，是因为零均值先验、相关结构与噪声共同决定了条件估计。

这个例子不是要学生相信所有函数都应服从该核。若真实响应有突变，过于平滑的核可能产生过度平滑甚至过度自信的结果。应改变核或噪声假设检查敏感性，并用独立数据检验预测覆盖，而不是仅因为程序输出了标准差就宣称不确定性完整。

\section{与本书其他方法的关系}
核SVM利用核定义相似性以形成判别函数，高斯过程则通过核定义函数值的协方差并形成条件分布。两者共享某些数学工具，目标与输出解释不同。随机森林树间离散、Bootstrap预测变化、GP条件方差也不能直接当成同一种不确定性。

高斯过程适合小到中等规模函数建模的教学，但标准稠密矩阵方法成本随样本增长较快。应用于光变曲线时，还要处理观测时刻、噪声与周期或长期趋势核；应用于物理响应时，要检查单位和领域约束。它是本书向概率建模延伸的入口，不是绕过数据与验证工作的捷径。
